\par\medskip
\noindent
5.5. To express conveniently the extreme factors in the filtrations
(5.1.10) and (5.1.11), in the case of semistable reduction, let us
introduce the twisted constant group schemes over $S$
\begin{equation}\tag{5.5.1}
 \underline M'_o=D(T_o),\qquad \underline M_o=D(T'_o),
\end{equation}
dual to the maximal tori $T_o,T'_o$ of the special fibres of $A,A'$
(SGA 3, X 5.7). Since $S$ is complete local, these group schemes may
be regarded as the special fibres of twisted constant schemes
\begin{equation}\tag{5.5.2}
 \underline M',\qquad \underline M\quad\hbox{over }S.
\end{equation}
Thus, if the torus $T_o$ is split (for example, if $k$ is separably
closed), that is, if $\underline M'_o$ is the constant group defined
by the ordinary group
\begin{equation}\tag{5.5.3}
 M'=\operatorname{Hom}_{k\text{-}\mathrm{gr}}(T_o,\mathbb G_{mk}),
\end{equation}
then $\underline M'$ is the constant group defined by $M'$:
\begin{equation*}
 \underline M'=M'_S.
\end{equation*}
One may moreover observe that, since $T_o$ and $T'_o$ are isogenous
(2.2.6), $T'_o$ is split whenever $T_o$ is. In this case one should
therefore also introduce
\begin{equation}\tag{5.5.4}
 M=\operatorname{Hom}_{k\text{-}\mathrm{gr}}(T'_o,\mathbb G_{mk}),
\end{equation}
and then
\begin{equation*}
 \underline M=M_S.
\end{equation*}
With this notation, for every integer $n$ there are evidently
isomorphisms
\begin{equation}\tag{5.5.5}
 \underline M'_n\simeq D(T_n),\qquad
 \underline M_n\simeq D(T'_n),
\end{equation}
where $T'_n$ is defined in terms of $A'$ as $T_n$ is in terms of $A$
(5.1.2). Hence there follow canonical isomorphisms (where
$\mathbb Z_\ell(1)=T_\ell(\mathbb G_{mS})$)
\begin{equation}\tag{5.5.6}
 T_\ell(T_n)\simeq
 \check{\underline M}'_n\otimes_{\mathbb Z}\mathbb Z_\ell(1),
 \qquad
 T_\ell(T'_n)\simeq
 \check{\underline M}_n\otimes_{\mathbb Z}\mathbb Z_\ell(1),
\end{equation}
and, on passing to the limit over $n$, canonical isomorphisms
\begin{equation}\tag{5.5.7}
 \left\{
 \begin{aligned}
  T_\ell(A^o)^t=T_\ell(\widehat T)
   &\simeq \check{\underline M}'\otimes\mathbb Z_\ell(1),\\[1ex]
  T_\ell(A^{\prime o})^t=T_\ell(\widehat T')
   &\simeq \check{\underline M}\otimes\mathbb Z_\ell(1).
 \end{aligned}
 \right.
\end{equation}
In the case of semistable reduction, duality, together with the
orthogonality theorem 5.2 and relation (5.2.2), then gives canonical
isomorphisms
\begin{equation}\tag{5.5.8}
 \left\{
 \begin{aligned}
  T_\ell(A_K)/T_\ell(A_K)^f
   &\simeq \underline M_K\otimes_{\mathbb Z}\mathbb Z_\ell,\\[1ex]
  T_\ell(A'_K)/T_\ell(A'_K)^f
   &\simeq \underline M'_K\otimes_{\mathbb Z}\mathbb Z_\ell.
 \end{aligned}
 \right.
\end{equation}
Likewise, the same formula (5.2.2) gives canonical isomorphisms
\begin{equation}\tag{5.5.9}
 \left\{
 \begin{aligned}
  T_\ell(A_K)/T_\ell(A_K)^t
   &\simeq D\bigl(T_\ell(A'_K)^f\bigr),\\[1ex]
  T_\ell(A'_K)/T_\ell(A'_K)^t
   &\simeq D\bigl(T_\ell(A_K)^f\bigr),
 \end{aligned}
 \right.
\end{equation}
where one should recall that the groups $T_\ell(A_K)^f$ and
$T_\ell(A'_K)^f$ occurring on the right are the restrictions to
$\eta$ of the pro-$\ell$ Barsotti--Tate groups $T_\ell(A^o)^f$ and
$T_\ell(A^{\prime o})^f$, respectively, defined over $S$. Here $D$
denotes duality in the sense of pro-Barsotti--Tate groups, namely
$\underline{\operatorname{Hom}}(-,\mathbb Z_\ell(1))$. Thus we have:

\medskip
\noindent
\underline{Proposition} 5.6. \underline{Suppose that $S$ is complete
and that $A_K$ has semistable reduction over $S$. Then, for every
prime number $\ell$ (which may equal the residue characteristic of
$S$), the pro-$\ell$ Barsotti--Tate group $T_\ell(A_K)^t$ is the
generic fibre of a pro-Barsotti--Tate group over $S$ ``of toric
type'', given by (5.5.7), whereas
$T_\ell(A_K)/T_\ell(A_K)^f$ is the generic fibre of a
pro-Barsotti--Tate group over $S$ of pro-\'{e}tale type, given by
(5.5.8). Finally, the groups $T_\ell(A_K)^f$ and
$T_\ell(A_K)/T_\ell(A_K)^t$ are the generic fibres of
pro-Barsotti--Tate groups defined over $S$, equal respectively to
$T_\ell(A^o)^f$ and to
$D(T_\ell(A^{\prime o})^f)=
\underline{\operatorname{Hom}}(T_\ell(A^{\prime o})^f,
\mathbb Z_\ell(1))$.}

\medskip
\noindent
5.6.1. Observe that, for an \'{e}tale covering $S'$ of $S$, splitting
$T_\ell(A_K)^t$ is equivalent to splitting
$T_\ell(A_K)/T_\ell(A_K)^f$: these conditions say respectively that
$S'$ splits the twisted constant groups $\underline M$ and
$\underline M'$, and these are isogenous because $T_o$ and $T'_o$
are. At the same time one sees that there is a smallest \'{e}tale
extension $S'$ of $S$ which does the job.

We also record the following easy consequence of the definitions in
2.2.3.
